\documentclass[conference]{IEEEtran}
\IEEEoverridecommandlockouts

\usepackage{cite}
\usepackage{amsmath,amssymb,amsfonts}
\usepackage{amsthm}
\usepackage{algorithm}
\usepackage{algpseudocode}
\usepackage{graphicx}
\usepackage{textcomp}
\usepackage{xcolor}
\usepackage{booktabs}
\usepackage{hyperref}
\usepackage{url}

\def\BibTeX{{\rm B\kern-.05em{\sc i\kern-.025em b}\kern-.08em
    T\kern-.1667em\lower.7ex\hbox{E}\kern-.125emX}}

\begin{document}

\title{Novelty Triggered Capacity Growth for Continual Learning\\
{\large Self Regulating Architecture Expansion via Gradient Trained Meta Parameters}}

\author{\IEEEauthorblockN{Susan Poudel}
\IEEEauthorblockA{\textit{Independent Researcher} \\
Fisk University, Nashville, TN\\
susanpdl77@gmail.com}
}

\maketitle

\begin{abstract}
Continual learning systems suffer from catastrophic forgetting when new tasks overwrite weights learned for earlier tasks. We present Novelty Triggered Capacity Growth (NCG). The method has two parts. First, three scalar meta parameters $\{\alpha, \beta, \lambda\}$ train with the network weights by gradient ascent on a Lagrangian meta loss, so exploration weight, complexity penalty, and regularisation strength need no hand tuned schedule. Second, a novelty signal from hidden layer activation entropy adds 64 hidden units only when the model has adapted, regularisation pressure crosses a threshold, and validation accuracy has plateaued. A gated knowledge embedding $K(t)$ updates across tasks as a compact summary alongside the weights. On Split-MNIST, forgetting drops by 21\% versus a parameter matched static baseline ($p = 0.012$, Welch $t$-test, $n = 10$ seeds). On Split-CIFAR-10, the drop is 64\% ($p < 0.0001$), and NCG beats EWC on that benchmark. Code: \url{https://github.com/Susanpdl/NCG}.
\end{abstract}

\begin{IEEEkeywords}
continual learning, catastrophic forgetting, capacity growth, meta parameters, gradient methods
\end{IEEEkeywords}

\section{Introduction}
Neural networks often fit one task well but lose accuracy when trained on a stream of tasks. Catastrophic forgetting \cite{mccloskey1989,french1999} happens because gradients for a new task move weights away from values that fit older tasks. That limits use when the input distribution shifts over time, for example in robotics, personalised assistants, medical monitoring, or fraud detection.

Most work falls into three groups. Regularisation methods such as EWC \cite{kirkpatrick2017} add a penalty so weights that mattered for past tasks do not move far. Replay methods such as GEM \cite{lopez2017} keep past samples and rehearse them. Parameter isolation methods such as PNN \cite{rusu2016} and DEN \cite{yoon2018} add capacity per task. Each line has limits: regularisation cuts plasticity, replay needs storage, and isolation often grows parameters with a fixed rule rather than a signal tied to need.

\textbf{Our contribution.} NCG is a parameter isolation method in which growth is not on a fixed schedule or task id. Growth follows from meta parameters trained with the weights to signal when more capacity helps. Specifically:
\begin{itemize}
    \item Three scalars $\alpha, \beta, \lambda$ update by gradient ascent on a Lagrangian meta loss instead of a fixed hyperparameter schedule.
    \item A normalised entropy novelty score and three conditions decide when to add units: the model has adapted, regularisation pressure is high, and validation accuracy has flattened.
    \item A gated knowledge embedding $K(t)$ stores a short exponential moving average of hidden states with a learned gate per dimension.
    \item We report ten seeds on Split-MNIST and Split-CIFAR-10 with ablations and baselines.
\end{itemize}

\section{Related Work}
\textbf{Regularisation methods.} EWC \cite{kirkpatrick2017} uses a diagonal Fisher approximation after each task and a quadratic penalty on important weights. SI \cite{zenke2017} tracks importance online. Both work best when tasks are similar; when tasks need different features, the plasticity stability trade off tightens.

\textbf{Replay and generative methods.} GEM \cite{lopez2017} and A-GEM \cite{chaudhry2019} store past minibatches and project gradients to limit interference. DGR \cite{shin2017} trains a generator for pseudo replay. Those approaches need buffers or a generator; NCG does not.

\textbf{Growing architectures.} DEN \cite{yoon2018} expands the net when a loss threshold is crossed. PNN \cite{rusu2016} adds a column per task with lateral links. PackNet \cite{mallya2018} prunes to free weights. HAT \cite{serra2018} uses task masks. In NCG, thresholds come from learned meta parameters instead of fixed rules.

\textbf{Meta learning for continual learning.} OML \cite{javed2019} and ANML \cite{beaulieu2020} use nested optimisation across tasks. NCG uses one training loop with no separate meta training stage.

\section{Method}
\subsection{Model State}
At any point in training, the NCG model is characterised by a state tuple:
\begin{equation}
    S(t) = \{W(t), K(t), \omega(t), d(t)\}
\end{equation}
where $W(t)$ are the network weights, $K(t) \in \mathbb{R}^{d(t)}$ is the knowledge embedding, $\omega(t) = \{\alpha(t), \beta(t), \lambda(t)\}$ are the meta parameters, and $d(t)$ is the current hidden dimension.

\subsection{Architecture}
\textbf{MLP backbone (Split-MNIST).}
\begin{equation}
    h = \mathrm{ReLU}(W_1 x + b_1 + K), \quad \hat{y} = W_2 h + b_2
\end{equation}
where $W_1 \in \mathbb{R}^{d \times n_{in}}$, $W_2 \in \mathbb{R}^{C \times d}$, and $K \in \mathbb{R}^d$ is added as a broadcast bias before the nonlinearity.

\textbf{CNN backbone (Split-CIFAR-10).} A two-layer CNN feature extractor (Conv2d $3 \to 32 \to 64$, kernel 3, ReLU, MaxPool $2 \times 2$, flattened to 1600-d) precedes the same FC structure.

\subsection{Knowledge Embedding with Gated Write}
The knowledge embedding $K(t) \in \mathbb{R}^{d(t)}$ stores a running summary of task-specific hidden representations. After each training batch, we update:
\begin{align}
    \bar{h} &= \frac{1}{B} \sum_{i=1}^{B} h_i \\
    g &= \sigma(W_g \bar{h} + b_g) \\
    K(t + 1) &= (1 - g) \odot K(t) + g \odot \bar{h}
\end{align}
where $W_g \in \mathbb{R}^{d \times d}$ is a learned gate projection, $\sigma$ is the sigmoid function, and $\odot$ denotes element-wise multiplication. The gate $g$ is trained jointly with $W_1, W_2$. Unlike replay buffers, $K$ is a learned summary vector, not raw stored inputs. When $d$ grows, the new $K$ dimensions are initialised to zero.

\subsection{Meta parameters}
Three scalar meta parameters govern the training objective. They are stored as unconstrained reals and mapped with smooth bijections:
\begin{align}
    \alpha &= \sigma(\tilde{\alpha}) \in (0, 1) \quad \text{exploration coefficient} \\
    \beta &= \mathrm{softplus}(\tilde{\beta}) \times 0.1 \in (0, \infty) \quad \text{complexity penalty} \\
    \lambda &= \sigma(\tilde{\lambda}) \in (0, 1) \quad \text{regularisation strength}
\end{align}
where $\sigma(x) = (1 + e^{-x})^{-1}$ and $\mathrm{softplus}(x) = \ln(1 + e^x)$.

\subsubsection{Training Loss}
At each step the network weights $W$ are updated by minimising:
\begin{equation}
    L_{\mathrm{train}} = L_{\mathrm{CE}} - \alpha^{\dagger} H(p) + (\beta^{\dagger} + \lambda^{\dagger}) \|W\|_F^2
\end{equation}
where $H(p) = -\sum_c p_c \log p_c$ is the predictive entropy of the current batch, $\|W\|_F^2 = \|W_1\|_F^2 + \|W_2\|_F^2$, and the superscript $\dagger$ denotes detached values (stop gradient). Weight updates use frozen meta values and meta updates use frozen weights so the two do not backprop into each other.

\subsubsection{Lagrangian Meta-Loss}
The meta parameters are updated by gradient ascent on a Lagrangian meta objective on a held out validation batch $V$. Define normalised norms:
\begin{align}
    \hat{A} &= \frac{\|W\|_F^2}{|W|} \\
    \hat{W} &= \sqrt{\hat{A} + \varepsilon}
\end{align}
where $|W|$ is the total number of weight parameters. The novelty signal on $V$ is:
\begin{equation}
    N_t = \frac{H(p_V)}{\ln C}
\end{equation}
where $C$ is the number of classes and $p_V = \mathrm{softmax}(f(V))$. The meta-loss is:
\begin{equation}
    L_{\mathrm{meta}} = L_{\mathrm{CE}}^V - \alpha(N_t - \tau_{\mathrm{nov}}) - \beta(\tau_A - \hat{A}) - \lambda(\tau_W - \hat{W})
\end{equation}
with thresholds $\tau_{\mathrm{nov}} = 0.3$, $\tau_A = 0.01$, $\tau_W = 0.1$. The update rule is gradient ascent:
\begin{equation}
    (\tilde{\alpha}, \tilde{\beta}, \tilde{\lambda}) \leftarrow (\tilde{\alpha}, \tilde{\beta}, \tilde{\lambda}) + \eta_m \nabla_{\tilde{\omega}} L_{\mathrm{meta}}
\end{equation}

The Lagrangian interpretation: the meta-loss is a constrained optimisation where $\alpha, \beta, \lambda$ are dual variables enforcing that (i) novelty stays above $\tau_{\mathrm{nov}}$, (ii) architecture norm stays below $\tau_A$, and (iii) weight norm stays below $\tau_W$. Gradient ascent on the duals is the standard Lagrangian update.

\subsection{Novelty Signal}
The novelty score $N_t \in [0, 1]$ is the normalised entropy of the predictive distribution on the current task's validation set:
\begin{itemize}
    \item $N_t \to 1$: near-uniform predictions; the task is novel, the model is uncertain.
    \item $N_t \to 0$: one class dominates; high confidence, task specialised.
    \item Task boundaries produce spikes, followed by monotonic decay as the model adapts.
\end{itemize}

\subsection{Growth Trigger}
Growth is triggered when all three conditions hold simultaneously:
\begin{align}
    C_1 &: N_t < \tau_n = 0.5 \\
    C_2 &: \lambda(t) > \tau_{\lambda} = 0.3 \\
    C_3 &: \max_{j \in [t-2, t]} A_j - \min_{j \in [t-2, t]} A_j < \tau_p = 0.005
\end{align}
where $A_j$ is the smoothed validation accuracy at epoch $j$ (three epoch moving average). A cooldown of $\Delta_{\min} = 5$ epochs prevents consecutive growth events.

$C_1$ requires adaptation before growth. $C_2$ requires strong enough regularisation to suggest the current width is tight. $C_3$ requires a flat validation curve so growth does not fire during fast improvement.

\subsection{Capacity Growth}
When the trigger fires, the hidden dimension expands from $d$ to $d + \Delta$ ($\Delta = 64$):
\begin{align}
    W_1' &= \begin{bmatrix} W_1 \\ U_1 \end{bmatrix}, \quad U_1 \sim \mathcal{N}(0, \sigma^2 I) \\
    W_2' &= \begin{bmatrix} W_2 & 0 \end{bmatrix}
\end{align}
where $\sigma = 0.1 \cdot \|W_1\|_F / \sqrt{|W_1|}$. Existing weights are copied; new $W_2$ columns start at zero so logits match the network immediately before widening. The knowledge buffer $K$ and gate layer $W_g$ are padded with zeros. The optimiser state is reset after growth. The maximum hidden width is $d_{\max} = 512$.

\subsection{Algorithm Summary}
\begin{algorithm}
\caption{NCG Training Loop}
\begin{algorithmic}[1]
\Require Tasks $\{T_k\}_{k=1}^{K}$, learning rates $\eta_w, \eta_m$
\State Initialise $W$, $K = 0$, $\omega = (0.5, 0.01, 0.5)$
\For{$k = 1, \ldots, K$}
    \For{epoch $= 1, \ldots, T$}
        \For{each minibatch $(x, y) \sim T_k$}
            \State Compute $h, \hat{y} = f(x; W, K)$
            \State $L \leftarrow L_{\mathrm{train}}(\omega^{\dagger})$
            \State $W \leftarrow W - \eta_w \nabla_W L$
            \State Update $K$ via gated write
        \EndFor
        \State Sample validation batch $V$ from $T_k$
        \State $L_m \leftarrow L_{\mathrm{meta}}$ on $V$
        \State $\tilde{\omega} \leftarrow \tilde{\omega} + \eta_m \nabla_{\tilde{\omega}} L_m$ \Comment{ascent}
        \If{$C_1 \wedge C_2 \wedge C_3$ and cooldown elapsed}
            \State Expand $W_1, W_2, K, W_g$
            \State Reinitialise optimiser state
        \EndIf
    \EndFor
\EndFor
\end{algorithmic}
\end{algorithm}

\subsection{Theoretical Analysis}
\textbf{Proposition 1 (Meta parameter gradient flow).} \textit{Let $L_{\mathrm{meta}}(\tilde{\omega})$ be twice continuously differentiable in $\tilde{\omega}$ with bounded Hessian. Under the gradient ascent update with sufficiently small $\eta_m$, any limit point $\tilde{\omega}^*$ satisfies $\nabla_{\tilde{\omega}} L_{\mathrm{meta}}(\tilde{\omega}^*) = 0$.}

\textit{Proof sketch.} The update $\tilde{\omega}_{t+1} = \tilde{\omega}_t + \eta_m g_t$ where $g_t = \nabla L_{\mathrm{meta}}$ is a standard gradient ascent iteration on a smooth objective. Applying the descent lemma to $-L_{\mathrm{meta}}$: if $\eta_m < 2/L$ where $L$ is the Lipschitz constant of $\nabla L_{\mathrm{meta}}$, then $\sum_t \|g_t\|^2 < \infty$, hence $\|g_t\|^2 \to 0$. Since $\omega$ is bounded (both $\sigma$ and $0.1 \cdot \mathrm{softplus}$ project to compact ranges), every trajectory has a convergent subsequence, and every limit point satisfies the stationarity condition. \qed

\textbf{Empirical stability test.} To test whether the limit point is stable, we perturb each meta parameter by $\delta = 0.1$ after training and run 20 ascent steps. The recovery ratio is:
\begin{equation}
    r = 1 - \frac{|\omega_{\mathrm{recovered}} - \omega^*|}{|\omega_{\mathrm{perturbed}} - \omega^*|}
\end{equation}
A value $r > 0.5$ constitutes evidence of a stable attractor. Results are reported in Section~\ref{sec:meta_dynamics}.

\section{Experimental Setup}
\subsection{Benchmarks}
\textbf{Split-MNIST.} MNIST \cite{lecun1998} is split into 5 binary classification tasks: (0/1), (2/3), (4/5), (6/7), (8/9). Standard 60k/10k train/test split filtered to two classes, with 15\% held-out validation. Inputs normalised to $[0, 1]$.

\textbf{Split-CIFAR-10.} CIFAR-10 \cite{krizhevsky2009} is split into 5 binary tasks: (airplane/automobile), (bird/cat), (deer/dog), (frog/horse), (ship/truck). Normalised with per-channel mean and standard deviation.

\subsection{Models}
\textbf{NCG.} Proposed model (MLP-256 initial hidden, growable to 512; CNN backbone for CIFAR-10).

\textbf{NCG-NoGrowth.} NCG with growth trigger disabled; meta parameters still updated.

\textbf{NCG-FixedMeta.} NCG with meta parameters frozen at $(\alpha, \beta, \lambda) = (0.5, 0.05, 0.73)$ (mean final values from NCG); growth trigger active.

\textbf{DEN.} Dynamically Expandable Networks \cite{yoon2018}, same backbone as NCG; grows when post-task validation loss $> 0.3$.

\textbf{StaticMLP-256 / 512.} Fixed-width 2-layer MLP trained with Adam, no continual learning mechanism.

\textbf{EWC.} Elastic Weight Consolidation \cite{kirkpatrick2017}, $\lambda_{\mathrm{EWC}} = 400$, diagonal Fisher computed on training data after each task.

\subsection{Training Details}
All NCG and ablation models: Adam optimiser, $\eta_w = 5 \times 10^{-4}$ (weights), $\eta_m = 10^{-2}$ (meta parameters), cosine annealing LR schedule per task, batch size 64, 10 epochs per task. StaticMLP and DEN: Adam, $\eta = 10^{-3}$. EWC: Adam, $\eta = 10^{-3}$. All experiments run over 10 seeds $\{42, 43, \ldots, 51\}$.

\subsection{Evaluation Metrics}
Following \cite{lopez2017}, after training on task $k$ we evaluate on all tasks $\{1, \ldots, K\}$. Let $a_{k,j}$ denote accuracy on task $j$ after training on task $k$.
\begin{align}
    \mathrm{AvgAcc} &= \frac{1}{K} \sum_{j=1}^{K} a_{K,j} \\
    \mathrm{Forgetting} &= \frac{1}{K-1} \sum_{j=1}^{K-1} \left[ \max_{k \leq K} a_{k,j} - a_{K,j} \right] \\
    \mathrm{BWT} &= \frac{1}{K-1} \sum_{j=1}^{K-1} (a_{K,j} - a_{j,j}) \\
    \mathrm{FWT} &= \frac{1}{K-1} \sum_{j=2}^{K} (a_{j-1,j} - b_j)
\end{align}
where $b_j = 0.5$ is the random-initialisation baseline for binary tasks.

\section{Results}
\subsection{Split-MNIST}
Table~\ref{tab:mnist} reports the full results. NCG achieves the lowest forgetting among all non-EWC methods (0.331 vs. 0.419 for StaticMLP-256). A Welch's $t$-test on per-seed forgetting scores confirms statistical significance: $t(10.8) = -3.06$, $p = 0.012$ (two-tailed), Cohen's $d = 1.47$ (large effect).

EWC outperforms all methods on this benchmark, consistent with known results: Split-MNIST is simple enough that Fisher-weighted regularisation succeeds without capacity growth.

\begin{table}[t]
\centering
\caption{Split-MNIST results. Mean $\pm$ std over 10 seeds. Lower forgetting is better. Bold: best forgetting among non-EWC methods.}
\label{tab:mnist}
\begin{tabular}{lcccc}
\toprule
Model & Avg Acc $\uparrow$ & Forgetting $\downarrow$ & BWT & FWT $\uparrow$ \\
\midrule
NCG (ours) & 0.551{\tiny$\pm$.016} & \textbf{0.331}{\tiny$\pm$.058} & $-$0.407{\tiny$\pm$.074} & 0.024{\tiny$\pm$.034} \\
NCG-NoGrowth & 0.552{\tiny$\pm$.005} & 0.373{\tiny$\pm$.024} & $-$0.466{\tiny$\pm$.043} & 0.039{\tiny$\pm$.037} \\
NCG-FixedMeta & 0.557{\tiny$\pm$.009} & 0.356{\tiny$\pm$.048} & $-$0.445{\tiny$\pm$.076} & 0.051{\tiny$\pm$.028} \\
DEN & 0.580{\tiny$\pm$.024} & 0.417{\tiny$\pm$.024} & $-$0.521{\tiny$\pm$.031} & 0.032{\tiny$\pm$.009} \\
StaticMLP-256 & 0.579{\tiny$\pm$.017} & 0.419{\tiny$\pm$.017} & $-$0.524{\tiny$\pm$.021} & 0.035{\tiny$\pm$.009} \\
StaticMLP-512 & 0.572{\tiny$\pm$.017} & 0.425{\tiny$\pm$.017} & $-$0.531{\tiny$\pm$.022} & 0.034{\tiny$\pm$.008} \\
EWC & 0.732{\tiny$\pm$.057} & 0.229{\tiny$\pm$.072} & $-$0.286{\tiny$\pm$.091} & 0.026{\tiny$\pm$.027} \\
\bottomrule
\end{tabular}
\end{table}

\textbf{Ablations.} NCG-NoGrowth (meta parameters on, no growth) lowers forgetting versus StaticMLP-256 (0.373 vs. 0.419), so regularisation from meta parameters helps. NCG-FixedMeta (growth on, frozen meta parameters) reaches 0.356, so growth also helps. Full NCG (0.331) is best; the two parts help together.

\subsection{Split-CIFAR-10}
Table~\ref{tab:cifar} gives Split-CIFAR-10 results. NCG reaches forgetting 0.083 versus 0.230 for StaticMLP-256, a 64\% relative reduction. Welch's $t$-test: $t(12.4) = -9.90$, $p < 0.0001$, Cohen's $d = 4.57$.

\begin{table}[t]
\centering
\caption{Split-CIFAR-10 results. Mean $\pm$ std over 10 seeds. Bold: best forgetting.}
\label{tab:cifar}
\begin{tabular}{lcccc}
\toprule
Model & Avg Acc $\uparrow$ & Forgetting $\downarrow$ & BWT & FWT $\uparrow$ \\
\midrule
NCG (ours) & 0.673{\tiny$\pm$.016} & \textbf{0.083}{\tiny$\pm$.044} & $-$0.086{\tiny$\pm$.057} & 0.061{\tiny$\pm$.009} \\
NCG-NoGrowth & 0.666{\tiny$\pm$.024} & 0.103{\tiny$\pm$.035} & $-$0.108{\tiny$\pm$.051} & 0.076{\tiny$\pm$.021} \\
NCG-FixedMeta & 0.673{\tiny$\pm$.020} & 0.096{\tiny$\pm$.028} & $-$0.119{\tiny$\pm$.040} & 0.077{\tiny$\pm$.019} \\
DEN & 0.688{\tiny$\pm$.006} & 0.222{\tiny$\pm$.005} & $-$0.278{\tiny$\pm$.007} & 0.088{\tiny$\pm$.011} \\
StaticMLP-256 & 0.683{\tiny$\pm$.011} & 0.230{\tiny$\pm$.010} & $-$0.288{\tiny$\pm$.013} & 0.086{\tiny$\pm$.007} \\
StaticMLP-512 & 0.687{\tiny$\pm$.013} & 0.227{\tiny$\pm$.014} & $-$0.284{\tiny$\pm$.017} & 0.088{\tiny$\pm$.009} \\
EWC & 0.702{\tiny$\pm$.016} & 0.163{\tiny$\pm$.023} & $-$0.203{\tiny$\pm$.029} & 0.088{\tiny$\pm$.012} \\
\bottomrule
\end{tabular}
\end{table}

On forgetting, NCG is below EWC (0.083 vs. 0.163), unlike on Split-MNIST. The CIFAR tasks differ more in appearance, so EWC's Fisher penalty tightens as tasks stack and can block new tasks. NCG adds width instead of freezing old weights.

DEN, the nearest expanding baseline, reaches 0.222 versus 0.083 for NCG. Both add units, but DEN uses a fixed loss cutoff while NCG requires all three trigger conditions.

\textbf{Ablations on CIFAR-10.} Growth alone (NCG-FixedMeta vs. StaticMLP-256) cuts forgetting by about 58\%. Meta parameters alone (NCG-NoGrowth) cut it by about 55\%. Full NCG at 0.083 beats both single part runs (0.096, 0.103), so the full model gains more than either part alone.

\section{Discussion}
\subsection{Novelty Signal Behaviour}
The novelty signal $N_t$ rises at task boundaries and falls within each task as the net fits the task. $N_t$ tracks shift in inputs without task ids or boundary labels.

\subsection{Growth Trigger Justification}
Growth tends to occur near task changes, after $N_t$ is low but while accuracy still moves. Across ten Split-MNIST seeds, NCG grows about 2.1 times on average (std 0.9), ending with $d \in \{256, 320, 384\}$, below $d_{\max} = 512$. On Split-CIFAR-10, growth fires about 3.2 times per run (std 0.7). Every run had at least one growth step.

\subsection{Meta parameter dynamics}
\label{sec:meta_dynamics}
All three meta parameters fall monotonically during training; the curves alone do not show whether that reflects intended adaptation or optimiser transients.

The perturbation test was run on trained model (seed 42), $\delta = 0.1$, 20 recovery steps:

\begin{table}[h]
\centering
\begin{tabular}{lccc}
\toprule
Parameter & $\omega^*$ & Recovery $r$ & Verdict \\
\midrule
$\alpha$ & 0.441 & 0.31 & Inconclusive \\
$\beta$ & 0.009 & 0.22 & Inconclusive \\
$\lambda$ & 0.635 & 0.29 & Inconclusive \\
\bottomrule
\end{tabular}
\end{table}

No parameter reaches $r > 0.5$. We therefore do not claim that $\alpha, \beta, \lambda$ converge to a stable fixed point. The proposition only gives a stationary point under smoothness; attractor stability is open. The reported 21\% and 64\% forgetting drops rest on the experiments, not on that stability claim.

\subsection{EWC Comparison}
EWC is stronger on Split-MNIST and weaker on Split-CIFAR-10. The Fisher penalty tightens as tasks stack and can limit updates on new tasks. NCG widens the net instead of penalising old weights as heavily, which helps when tasks look different. Where the two methods swap rank will depend on task similarity and difficulty.

\subsection{Forward Transfer}
On Split-MNIST, NCG has the lowest FWT (0.024), below NCG-NoGrowth (0.039) and NCG-FixedMeta (0.051). Growth resets the optimiser state, which may briefly hurt features that would help the next task. A practical next step is to keep optimiser state for old units and only initialise state for new units.

\subsection{Computational Overhead}
NCG's primary overhead over a static baseline is the meta-update step: one additional forward pass on the validation batch and a backward pass through three scalars. This adds approximately 8\% wall-clock time per epoch on CPU. The gated write is $O(d^2)$ per batch; at $d = 320$, this is negligible. The growth operation is $O(d \cdot \Delta)$ and occurs at most $\lfloor (d_{\max} - d_0)/\Delta \rfloor = 4$ times per run.

\section{Limitations}
\textbf{Benchmark scope.} We only use Split-MNIST and Split-CIFAR-10 with two classes per task. We did not run Split-CIFAR-100 or permuted MNIST here.

\textbf{Meta parameter convergence.} The perturbation test does not show a stable fixed point. The theory section only argues for a stationary point, not that it attracts nearby points. The tables do not rely on that claim.

\textbf{Thresholds.} $\tau_n$, $\tau_{\lambda}$, and $\tau_p$ were set from a small set of trial runs, not an exhaustive grid.

\textbf{Forward transfer.} Full NCG has lower FWT than the ablations on Split-MNIST; we do not have a full explanation.

\textbf{Architectures.} We use MLPs and a small CNN. The public code includes a growth hook for transformer FFN blocks, which we did not benchmark.

\section{Conclusion}
We described NCG, which pairs gradient trained meta parameters with a novelty based growth rule to reduce forgetting in continual learning. On Split-MNIST, forgetting falls by 21\% versus a parameter matched static baseline ($p = 0.012$, $n = 10$ seeds). On Split-CIFAR-10, the drop is 64\% ($p < 0.0001$), and NCG beats EWC there, which fits tasks that need different visual features.

Ablations on CIFAR-10 show both parts matter: meta parameters alone cut forgetting by about 55\%, growth alone by about 58\%, and the full model reaches 64\%. Meta parameters change when growth fires, not only whether the net grows.

The gated vector $K$ gives a small task summary without storing raw past data, which differs from many expansion baselines. Left for later work: the FWT drop under full NCG, tighter theory on meta parameter limits, and tests on harder splits such as Split-CIFAR-100 or permuted MNIST.

Code, data, and checkpoints: \url{https://github.com/Susanpdl/NCG}.

\begin{thebibliography}{15}
\bibitem{mccloskey1989} M. McCloskey and N. J. Cohen, ``Catastrophic interference in connectionist networks: The sequential learning problem,'' in \textit{Psychology of Learning and Motivation}, vol.~24, pp.~109--165, 1989.

\bibitem{french1999} R. M. French, ``Catastrophic forgetting in connectionist networks,'' \textit{Trends in Cognitive Sciences}, vol.~3, no.~4, pp.~128--135, 1999.

\bibitem{kirkpatrick2017} J. Kirkpatrick \textit{et al.}, ``Overcoming catastrophic forgetting in neural networks,'' \textit{Proc. National Academy of Sciences}, vol.~114, no.~13, pp.~3521--3526, 2017.

\bibitem{lopez2017} D. L\'opez-Paz and M. A. Ranzato, ``Gradient episodic memory for continual learning,'' in \textit{NeurIPS}, 2017.

\bibitem{rusu2016} A. A. Rusu \textit{et al.}, ``Progressive neural networks,'' \textit{arXiv:1606.04671}, 2016.

\bibitem{yoon2018} J. Yoon, E. Yang, J. Lee, and S. J. Hwang, ``Lifelong learning with dynamically expandable networks,'' in \textit{ICLR}, 2018.

\bibitem{zenke2017} F. Zenke, B. Poole, and S. Ganguli, ``Continual learning through synaptic intelligence,'' in \textit{ICML}, 2017.

\bibitem{shin2017} H. Shin, J. K. Lee, J. Kim, and J. Kim, ``Continual learning with deep generative replay,'' in \textit{NeurIPS}, 2017.

\bibitem{chaudhry2019} A. Chaudhry \textit{et al.}, ``Efficient lifelong learning with A-GEM,'' in \textit{ICLR}, 2019.

\bibitem{mallya2018} A. Mallya and S. Lazebnik, ``PackNet: Adding multiple tasks to a single network by iterative pruning,'' in \textit{CVPR}, 2018.

\bibitem{serra2018} J. Serra \textit{et al.}, ``Overcoming catastrophic forgetting with hard attention to the task,'' in \textit{ICML}, 2018.

\bibitem{javed2019} K. Javed and M. White, ``Meta-learning representations for continual learning,'' in \textit{NeurIPS}, 2019.

\bibitem{beaulieu2020} S. Beaulieu \textit{et al.}, ``Learning to continually learn,'' in \textit{ECAI}, 2020.

\bibitem{lecun1998} Y. LeCun, L. Bottou, Y. Bengio, and P. Haffner, ``Gradient-based learning applied to document recognition,'' \textit{Proc. IEEE}, vol.~86, no.~11, pp.~2278--2324, 1998.

\bibitem{krizhevsky2009} A. Krizhevsky, ``Learning multiple layers of features from tiny images,'' Tech. Rep., University of Toronto, 2009.
\end{thebibliography}

\end{document}
